% !TeX root = ../教材题库_全部条目_转写工作稿.tex
% Source: XJ-B1-C02-EX-P034-Y0712
% 第2章 / 2.1 相等关系与不等关系 / 教材 p.34 / PDF 第41页
% OCR未校对：公式与文字均可能有误；未完成数学TeX转写；禁止直接发给学生
% 原题及相邻原解截图（核对用）：../题目原图/XJ-B1-C02-EX-P034-Y0712-01.webp
% 本文件只在转写完成后作为题干使用；OCR 条目尚不是可靠题干。
\begin{bankitem}{2.1 例题 例题2}{第2章 / 2.1 相等关系与不等关系 / 教材 p.34 / PDF 第41页}{XJ-B1-C02-EX-P034-Y0712}{OCR未校对：公式与文字均可能有误；未完成数学TeX转写；禁止直接发给学生}
\begin{ocrdraft}
[教材p.34 / PDF第41页]
\par
00g糖水中含有bg糖，若再添加g糖（其中a>b>0，m>0)，生活常
\par
识告诉我们：添加的糖完全溶解后，糖水会更甜．根据这个生活常识，你能提炼出
\par
一个不等式吗？试给出证明．
\par
解因为加糖后糖水更甜，即糖水中糖的浓度变大，所以提炼出的不等式为：
\par
b十mb
\par
>一，其中a>b>O,m>O.
\par
下面用作差比较法给出证明．
\par
假设有一种机器可以
\par
乃十m乃一鼠乃十m）一乃（a十m）一m一b）
\par
抽取糖水中的糖，生活常
\par
十m）
\par
a(a十m)
\par
识告诉我们：若从糖水中
\par
因为0乃，m都是正数，且a>b,
\par
抽掉mg糖，则糖水会变
\par
所以
\par
a十m>O，磊一乃>0．
\par
淡，于是提炼出一个不等
\par
m(a—b)
\par
式：若口>b>m>0，则
\par
因此
\par
鼠十”D
\par
<一．你能证明这
\par
b十mb
\par
个不等式吗？
\end{ocrdraft}
\end{bankitem}
